% GENERATED FILE. Edit canonical lessons, then run npm run generate:books.
% canonical-source-hash: fnv1a64:251c95b07fb3023a
% canonical-lessons: TE-C253-civara, TE-C253-eppudaina, TE-C253-prastutam, TE-C253-purvam, TE-C253-tellavarujamu

\chapter{End, Ever, At present, In the past, Dawn}
\label{ch:te-end-ever-at-present-in-the-past-dawn}
\hlchaptermodality{253}

\begin{chapteropening}
\textbf{By the end of this chapter:} \emph{I can say the end, ever, at present, in the past and dawn.}

Complete the last lesson of chapter 253: I can say the end, ever, at present, in the past and dawn.
\end{chapteropening}

\section[\texorpdfstring{\te{చివర} (civara)}{చివర (civara)}]{\te{చివర} (civara) — the end}

\label{lesson:TE-C253-civara}



\begin{quote}
Before the new one: say the Telugu for Vinayaka Chavithi, then the Telugu for Bathukamma.
\end{quote}

\subsection*{You'll want to know: \te{చివర}}
\textbf{\te{చివర}} — \emph{civara} — \textquotedblleft{}the end\textquotedblright{}.

\begin{grammarlens}[title={What you've built}]
One more word for this chapter.
\end{grammarlens}

\subsection*{Your turn}
\begin{itemize}
\raggedright
  \item \emph{Say it:} \emph{civara}
  \item \emph{Say it:} \emph{civara}, once more
  \item \emph{Say it:} say \emph{civara}
  \item \emph{Recall:} say the Telugu for Vinayaka Chavithi, then the Telugu for Bathukamma, then say \emph{civara} again
\end{itemize}

\begin{tcolorbox}[breakable,colback=teal!4,colframe=teal!35!black,title={Before you move on}]
What does \te{చివర} mean? (\textquotedblleft{}The end\textquotedblright{}.) Say it once more.
\end{tcolorbox}
\section[\texorpdfstring{\te{ఎప్పుడైనా} (eppuḍainā)}{ఎప్పుడైనా (eppuḍainā)}]{\te{ఎప్పుడైనా} (eppuḍainā) — ever, at any time}

\label{lesson:TE-C253-eppudaina}



\begin{quote}
Before the new one: say the Telugu for Bathukamma, then the Telugu for the end.
\end{quote}

\subsection*{You'll want to know: \te{ఎప్పుడైనా}}
\textbf{\te{ఎప్పుడైనా}} — \emph{eppuḍainā} — \textquotedblleft{}ever, at any time\textquotedblright{}.

\begin{grammarlens}[title={What you've built}]
One more word for this chapter.
\end{grammarlens}

\subsection*{Your turn}
\begin{itemize}
\raggedright
  \item \emph{Say it:} \emph{eppuḍainā}
  \item \emph{Say it:} \emph{eppuḍainā}, once more
  \item \emph{Say it:} say \emph{eppuḍainā}
  \item \emph{Recall:} say the Telugu for Bathukamma, then the Telugu for the end, then say \emph{eppuḍainā} again
\end{itemize}

\begin{tcolorbox}[breakable,colback=teal!4,colframe=teal!35!black,title={Before you move on}]
What does \te{ఎప్పుడైనా} mean? (\textquotedblleft{}Ever, at any time\textquotedblright{}.) Say it once more.
\end{tcolorbox}
\section[\texorpdfstring{\te{ప్రస్తుతం} (prastutaṁ)}{ప్రస్తుతం (prastutaṁ)}]{\te{ప్రస్తుతం} (prastutaṁ) — at present, currently}

\label{lesson:TE-C253-prastutam}



\begin{quote}
Before the new one: say the Telugu for the end, then the Telugu for ever.
\end{quote}

\subsection*{You'll want to know: \te{ప్రస్తుతం}}
\textbf{\te{ప్రస్తుతం}} — \emph{prastutaṁ} — \textquotedblleft{}at present, currently\textquotedblright{}.

\begin{grammarlens}[title={What you've built}]
One more word for this chapter.
\end{grammarlens}

\subsection*{Your turn}
\begin{itemize}
\raggedright
  \item \emph{Say it:} \emph{prastutaṁ}
  \item \emph{Say it:} \emph{prastutaṁ}, once more
  \item \emph{Say it:} say \emph{prastutaṁ}
  \item \emph{Recall:} say the Telugu for the end, then the Telugu for ever, then say \emph{prastutaṁ} again
\end{itemize}

\begin{tcolorbox}[breakable,colback=teal!4,colframe=teal!35!black,title={Before you move on}]
What does \te{ప్రస్తుతం} mean? (\textquotedblleft{}At present, currently\textquotedblright{}.) Say it once more.
\end{tcolorbox}
\section[\texorpdfstring{\te{పూర్వం} (pūrvaṁ)}{పూర్వం (pūrvaṁ)}]{\te{పూర్వం} (pūrvaṁ) — in the past, formerly}

\label{lesson:TE-C253-purvam}



\begin{quote}
Before the new one: say the Telugu for ever, then the Telugu for at present.
\end{quote}

\subsection*{You'll want to know: \te{పూర్వం}}
\textbf{\te{పూర్వం}} — \emph{pūrvaṁ} — \textquotedblleft{}in the past, formerly\textquotedblright{}.

\begin{grammarlens}[title={What you've built}]
One more word for this chapter.
\end{grammarlens}

\subsection*{Your turn}
\begin{itemize}
\raggedright
  \item \emph{Say it:} \emph{pūrvaṁ}
  \item \emph{Say it:} \emph{pūrvaṁ}, once more
  \item \emph{Say it:} say \emph{pūrvaṁ}
  \item \emph{Recall:} say the Telugu for ever, then the Telugu for at present, then say \emph{pūrvaṁ} again
\end{itemize}

\begin{tcolorbox}[breakable,colback=teal!4,colframe=teal!35!black,title={Before you move on}]
What does \te{పూర్వం} mean? (\textquotedblleft{}In the past, formerly\textquotedblright{}.) Say it once more.
\end{tcolorbox}
\section[\texorpdfstring{\te{తెల్లవారుజాము} (tellavārujāmu)}{తెల్లవారుజాము (tellavārujāmu)}]{\te{తెల్లవారుజాము} (tellavārujāmu) — dawn, the early hours}

\label{lesson:TE-C253-tellavarujamu}



\begin{quote}
Before the new one: say the Telugu for at present, then the Telugu for in the past.
\end{quote}

\subsection*{You'll want to know: \te{తెల్లవారుజాము}}
\textbf{\te{తెల్లవారుజాము}} — \emph{tellavārujāmu} — \textquotedblleft{}dawn, the early hours\textquotedblright{}.

\begin{grammarlens}[title={What you've built}]
One more word for this chapter.
\end{grammarlens}

\subsection*{Your turn}
\begin{itemize}
\raggedright
  \item \emph{Say it:} \emph{tellavārujāmu}
  \item \emph{Say it:} \emph{tellavārujāmu}, once more
  \item \emph{Say it:} say \emph{tellavārujāmu}
  \item \emph{Recall:} say the Telugu for at present, then the Telugu for in the past, then say \emph{tellavārujāmu} again
\end{itemize}

\begin{tcolorbox}[breakable,colback=teal!4,colframe=teal!35!black,title={Before you move on}]
What does \te{తెల్లవారుజాము} mean? (\textquotedblleft{}Dawn, the early hours\textquotedblright{}.) Say it once more.
\end{tcolorbox}
